\documentclass{article}
\usepackage{amsmath,hyperref}
\title{A zone's old idle sample can decide another queue's service}
\begin{document}
\maketitle

Two replays have the same free spare slots and the same requests waiting for them. One previous visit was seen by the zone timer in just one replay. That changes which spare is returned first, even though the visitor is gone in both. Forwarding the real release messages into the reservation allocator gives one granted move in one replay and six in the other. Under a separate physical-clearance guard the six moves finish on the latter replay, but five of them are still physically pending when their grants are issued.

\paragraph{One history through both handlers.}
Use the selected \href{https://github.com/open-rmf/rmf_ros2/tree/16a3457196c6edf2f08a1ba4f497bf891404cbb7}{Open-RMF sources} from the earlier \href{https://github.com/eternal-november/eos-research-programmes/pull/451}{zone sweep} and \href{https://github.com/eternal-november/eos-research-programmes/pull/452}{release-order} studies. This replay connects their compiled portions. The ordinary allocator generates the same ticket identities that the zone manager temporarily holds; the zone release publisher's ticket goes back to that same allocator on the same clock. No comparison splices the release from one run onto the reservations of another run. Delivery is in-process without extra lag. Initial navigation, setup, hardware occupancy, and callback times are fixture inputs; no ROS scheduler or robot motion stack is executed.

Let the idle timeout be \(C\ge3\) ticks. Two zones hold one spare apiece, \(A\) and \(B\). The first idle check of \(A\) is at 0. The hold for \(B\) is installed later and its first idle check is at 1. A temporary visitor is granted \(A\). In one trace the visitor departs before the check at 1. In the other the check sees the visitor en route and clears the timestamp, then the visitor departs. By the check at 2 the visitor is gone and both holdings are identical; the pending requests, existing allocations and any release publications agree. Only the sampled timer history differs. The code returns the spares in order \(A@C,B@(C+1)\) or \(B@(C+1),A@(C+2)\). There is no dependence on breaking a tie between simultaneous callback publications in this test.

An earlier flexible request \(F\) accepts either spare. A constrained request \(Q\) accepts only \(A\) and currently occupies \(P_0\). Follower \(i\) waits at \(P_i\) and needs \(P_{i-1}\). If \(A\) is returned first, the native queue assigns it to \(F\); freeing \(B\) cannot admit \(Q\), so none of the followers receives a grant. Reversing the release order assigns \(F\) to \(B\) and then \(Q\) to \(A\). The latter grant frees the logical waiting ticket at \(P_0\), and induction through the wait chain yields \(n+2\) grants for \(n\) followers. A control placing \(Q\) first on the \(A\) queue admits all jobs in either timer history. Both handlers can satisfy their local timer and allocation rules throughout. An outer guarantee of service for the whole queue needs the joint timing and contention assumption.

\paragraph{Actual effect is a separate cut.}
Add a source-independent guard for the declared physical fixture: only move into a destination after its previous occupant has cleared, reserve it during motion, and keep the old waiting position occupied until that robot completes its move. Assume no other conflicts or path constraints, unit-capacity positions, and a common movement duration \(d>0\). The grants along the chain arrive at the same logical tick; they do not clear the waiting positions. The first constrained motion can finish at \(t_A+d\), and follower \(i\) at \(t_A+(i+1)d\) for \(1\le i\le n\). These times are achieved by starting each as soon as the preceding physical clearance occurs; no faster safe schedule is possible under this movement model. In the branch serving everyone, \(F\) uses \(B\) and finishes separately at \(t_B+d\). The exact count by deadline \(D\) is
\[
  \mathbf 1_{D\ge t_B+d}+
  \max\!\left(0,\min\!\left(n+1,\left\lfloor \frac{D-t_A}{d}\right\rfloor\right)\right).
\]
The one-grant branch gives just \(\mathbf 1_{D\ge t_A+d}\). For \(C=30\), \(d=2\), and four followers, the two joined source traces give one guarded completion by 32 and six by 42. Before accepting all six by a deadline, later consumers need the actual release time as well as the queue order; a count of logical grants is insufficient. Delaying a release can improve total completion in this family while worsening completion of a particular urgent job.

\paragraph{Comparison and limits.}
The exact factored direct model first computes the timer order and then the queue consequences. It agrees with the joined native execution in all 2,940 cases. The result shows which information must cross this particular interface: named live tickets, the ordering of their releases when alternative destinations contend, and clearance times for downstream physical admission. A receiver seeing only the current holding list cannot guarantee the better branch across these histories. If a different allocator has lawful access to both physically clear destinations before assigning either, it can instead plan the better assignment directly; that evidence and any buffering time must be supplied and charged. This calibration is neither a new allocator nor an independent mission test.
\end{document}
